\begingroup
\interfootnotelinepenalty=10000
\Needspace{12\baselineskip}
\begin{center}\textit{Ein ähnliches Verfahren zum Beweis des Satzes aus Artikel 114.}\end{center}

\Article{145}
In den vorigen Beweisen haben wir für $e$ stets einen geraden Wert genommen (Artikel 137\,..\,144). Man hätte auch einen ungeraden Wert verwenden können, hätte dann aber weitere Unterscheidungen einführen müssen. Wer an diesen Untersuchungen Freude hat, wird seine Kräfte mit Nutzen daran erproben, diese Fälle auszuarbeiten. Außerdem hätten dann die Sätze über die Reste $+2$ und $-2$ vorausgesetzt werden müssen. Da unser Beweis aber ohne diese Sätze auskommt, gewinnen wir ein neues Verfahren, um sie zu beweisen. Das ist durchaus beachtenswert, denn die oben verwendeten Methoden zum Beweis, dass $\pm2$ ein Rest jeder Primzahl der Form $8n+1$ ist, mögen weniger direkt erscheinen. Wir setzen voraus, dass die übrigen Fälle (zu den Primzahlen der Formen $8n+3$, $8n+5$, $8n+7$) mit den oben angegebenen Methoden bewiesen sind und jener Satz nur anhand von Beispielen gefunden wurde. Diese Beobachtung werden wir durch die folgenden Überlegungen zur Gewissheit erheben.

Wäre $\pm2$ nicht ein Rest aller Primzahlen der Form $8n+1$, so sei die kleinste Primzahl dieser Form, von der $\pm2$ ein Nichtrest ist, $=a$. Für alle Primzahlen unterhalb von $a$ gilt dann der Satz. Man nehme eine Primzahl $<\frac{1}{2}a$, von der $a$ ein Nichtrest ist (dass es eine solche gibt, folgt leicht aus Artikel 129). Diese sei $=p$. Nach dem Fundamentalsatz gilt dann $pNa$, also $\pm2pRa$. --- Sei daher $e^2\equiv2p\ (\mathrm{mod.}\ a)$, wobei $e$ ungerade und $<a$ ist. Nun sind zwei Fälle zu unterscheiden.

I. Ist $e$ nicht durch $p$ teilbar, so sei $e^2=2p+aq$. Dann ist $q$ positiv, von der Form $8n+7$ oder $8n+3$ (je nachdem, ob $p$ die Form $4n+1$ oder $4n+3$ hat), $<a$ und nicht durch $p$ teilbar. Man teile alle Primfaktoren von $q$ in vier Klassen ein: Es seien $e$ von der Form $8n+1$, $f$ von der Form $8n+3$, $g$ von der Form $8n+5$ und $h$ von der Form $8n+7$. Das Produkt der Faktoren der ersten Klasse sei $E$, die Produkte der Faktoren der zweiten, dritten und vierten Klasse entsprechend $F,G.H$\HistoricalNote{*)}{Enthält eine Klasse keine Faktoren, so muss man für deren Produkt $1$ schreiben.}.

\SourceLeaf{116}Betrachten wir nun \textit{erstens} den Fall, dass $p$ die Form $4n+1$ oder $q$ die Form $8n+7$ hat. Dann gilt offensichtlich $2RE$, $2RH$, also $pRE$, $pRH$ und schließlich $ERp$, $HRp$. Ferner ist $2$ ein Nichtrest jedes Faktors der Form $8n+3$ oder $8n+5$; dasselbe gilt für $p$. Daher ist jeder solche Faktor ein Nichtrest von $p$. Es folgt leicht, dass $FG$ ein Rest von $p$ ist, wenn $f+g$ gerade ist, und ein Nichtrest, wenn $f+g$ ungerade ist. Aber $f+g$ kann nicht ungerade sein: Durch Aufzählen aller Fälle sieht man leicht, dass $EFGH$ oder $q$ dann die Form $8n+3$ oder $8n+5$ hätte, wenn $f+g$ ungerade wäre, welche Werte die einzelnen $e,f,g,h$ auch hätten. Das widerspricht der Voraussetzung. Also gilt $FGRp$, $EFGHRp$ oder $qRp$ und schließlich wegen $aqRp$ auch $aRp$, entgegen der Voraussetzung. Hat \textit{zweitens} $p$ die Form $4n+3$, so kann man ähnlich zeigen, dass $pRE$, also $ERp$, und $-pRF$, also $FRp$, gilt; schließlich ist $g+h$ gerade, also gilt $GHRp$, woraus zuletzt $qRp$, $aRp$ folgt, entgegen der Voraussetzung.

II. Ist $e$ durch $p$ teilbar, so lässt sich der Beweis auf ähnliche Weise führen und von geübten Lesern (für die allein dieser Artikel geschrieben ist) ohne Schwierigkeit ausarbeiten. Der Kürze wegen lassen wir ihn weg.

\Needspace{10\baselineskip}
\begin{center}\textit{Lösung des allgemeinen Problems.}\end{center}

\Article{146}
Mit dem Fundamentalsatz und den Aussagen über die Reste $-1$ und $\pm2$ lässt sich immer bestimmen, ob eine beliebige gegebene Zahl ein Rest oder Nichtrest einer gegebenen Primzahl ist. Es wird aber nützlich sein, auch die anderen oben angegebenen Ergebnisse hier noch einmal zusammenzustellen, damit alles zur Lösung des folgenden Problems Nötige beisammensteht.

\textsc{Aufgabe:} \textit{Sind zwei beliebige Zahlen $P$, $Q$ gegeben, so bestimme man, ob die zweite, $Q$, ein Rest oder Nichtrest der ersten, $P$, ist.}

\textit{Lösung.} I. Sei $P=a^{\alpha}b^{\beta}c^{\gamma}$ usw., wobei $a,b,c$ usw. verschiedene positive Primzahlen bezeichnen (denn $P$ ist offensichtlich im Betrag zu nehmen). Der Kürze wegen nennen wir in diesem Artikel einfach die \textit{Beziehung} zweier Zahlen $x,y$ diejenige, nach der die erste, $x$, ein Rest oder Nichtrest der zweiten, $y$, ist. Die Beziehung von $Q,P$ hängt also von den Beziehungen von $Q,a^{\alpha}$; $Q,b^{\beta}$ usw. ab (Artikel 105).

II. Um die Beziehung von $Q,a^{\alpha}$ zu bestimmen (für die übrigen $Q,b^{\beta}$ usw. gilt dasselbe), sind zwei Fälle zu unterscheiden.

\SourceLeaf{117}1. Ist $Q$ durch $a$ teilbar, so setze man $Q=Q'a^e$, wobei $Q'$ nicht durch $a$ teilbar ist. Ist dann $e=\alpha$ oder $e>\alpha$, so gilt $QRa^{\alpha}$; ist dagegen $e<\alpha$ und der Exponent ungerade, so gilt $QNa^{\alpha}$; ist schließlich $e<\alpha$ und der Exponent gerade, so hat $Q$ zu $a^{\alpha}$ dieselbe Beziehung wie $Q'$ zu $a^{\alpha-e}$. Dieser Fall ist damit zurückgeführt auf

2. den Fall, dass $Q$ nicht durch $a$ teilbar ist. Hier unterscheiden wir wiederum zwei Fälle.

$(A)$ Wenn $a=2$. Dann gilt stets $QRa^{\alpha}$, sobald $\alpha=1$ ist. Ist dagegen $\alpha=2$, so muss $Q$ die Form $4n+1$ haben; ist schließlich $\alpha=3$ oder $>3$, so muss $Q$ die Form $8n+1$ haben. Ist diese Bedingung erfüllt, so gilt $QRa^{\alpha}$.

$(B)$ Wenn $a$ eine andere Primzahl ist. Dann hat $Q$ zu $a^{\alpha}$ dieselbe Beziehung wie zu $a$ (siehe Artikel 101).

III. Die Beziehung einer beliebigen Zahl $Q$ zu einer ungeraden Primzahl $a$ wird so ermittelt. Ist $Q>a$, so ersetze man $Q$ durch seinen kleinsten positiven Rest nach dem Modul $a$\HistoricalNote{*)}{\textit{Rest} im Sinne von Artikel 4. --- Meist ist es besser, den im \textit{Betrag} kleinsten Rest zu nehmen.}. Dieser hat zu $a$ dieselbe Beziehung wie $Q$.

Ferner zerlege man $Q$ oder die dafür genommene Zahl in ihre Primfaktoren $p,p',p''$ usw.; der Faktor $-1$ ist hinzuzufügen, wenn $Q$ negativ ist. Bekanntlich hängt dann die Beziehung von $Q$ zu $a$ von den Beziehungen der einzelnen $p,p',p''$ usw. zu $a$ ab. Sind nämlich unter jenen Faktoren $2m$ Nichtreste von $a$, so gilt $QRa$; sind es $2m+1$, so gilt $QNa$. Kommen unter den Faktoren $p,p',p''$ usw. zwei oder vier oder sechs oder allgemein $2k$ gleiche vor, so kann man sie offensichtlich ohne Bedenken weglassen.

IV. Finden sich unter den Faktoren $p.p',p''$ die Zahlen $-1$ und $2$, so lässt sich ihre Beziehung zu $a$ aus den Artikeln 108, 112, 113, 114 bestimmen. Die Beziehung der übrigen zu $a$ hängt von der Beziehung von $a$ zu ihnen ab (Fundamentalsatz und Aussagen aus Artikel 131). Sei $p$ einer dieser Faktoren. Behandelt man die Zahlen $a,p$ genauso wie zuvor $Q$ und $a$, die entsprechend größer sind, so findet man, dass sich die Beziehung von $a$ zu $p$ entweder nach den Artikeln 108---114 bestimmen lässt (wenn nämlich der kleinste Rest von $a$ $(\mathrm{mod.}\ p)$ keine ungeraden Primfaktoren hat) oder weiter von der Beziehung von $p$ zu gewissen Primzahlen unterhalb von $p$ abhängt. Dasselbe gilt für die übrigen Faktoren $p',p''$ usw. Nun sieht man leicht, dass \SourceLeaf{118}die Fortsetzung dieses Verfahrens schließlich zu Zahlen führt, deren Beziehungen nach den Aussagen aus den Artikeln 108---114 bestimmt werden können. Ein Beispiel macht dies deutlicher.

\textit{Beispiel.} Gesucht ist die Beziehung von $+453$ zu $1236$. Es gilt $1236=4.3.103$; $+453R4$ nach II. 2 $(A)$; $+453R3$ nach II. 1. Es bleibt also die Beziehung von $+453$ zu $103$ zu untersuchen. Sie ist dieselbe wie die von $+41$ $(\equiv453,\ \mathrm{mod.}\ 103)$ zu $103$, dieselbe wie die von $+103$ zu $41$ (nach dem Fundamentalsatz) oder die von $-20$ zu $41$. Es gilt aber $-20R41$, denn $-20=-1.2.2.5$; $-1R41$ (Artikel 108) und $+5R41$, weil $41\equiv1$ und die Zahl somit ein Rest von $5$ ist (nach dem Fundamentalsatz). Daraus folgt $+453R103$ und schließlich $+453R1236$. Tatsächlich gilt $453\equiv297^2\ (\mathrm{mod.}\ 1236)$.

\Needspace{12\baselineskip}
\begin{center}\textit{Lineare Formen, die alle Primzahlen enthalten, von denen eine beliebige gegebene Zahl ein Rest oder Nichtrest ist.}\end{center}

\Article{147}
Ist eine beliebige Zahl $A$ gegeben, so lassen sich bestimmte \textit{Formeln} angeben, die alle zu $A$ teilerfremden Zahlen enthalten, von denen $A$ ein Rest ist, oder alle Zahlen, die \textit{Teiler} von Zahlen der Form $xx-A$ sein können (wobei $xx$ ein unbestimmtes Quadrat bezeichnet)\HistoricalNote{*)}{Solche Zahlen nennen wir einfach \textit{Teiler von $xx-A$}; daraus ist klar, was \textit{Nichtteiler} sind.}. Der Kürze wegen betrachten wir aber nur die Teiler, die ungerade und zu $A$ teilerfremd sind, da sich die übrigen Fälle leicht darauf zurückführen lassen.

Zunächst sei $A$ entweder eine positive Primzahl der Form $4n+1$ oder eine negative der Form $4n-1$. Nach dem Fundamentalsatz sind dann alle Primzahlen, die, positiv genommen, Reste von $A$ sind, Teiler von $xx-A$; alle Primzahlen dagegen (außer der Zahl $2$, die immer ein Teiler ist), die Nichtreste von $A$ sind, sind Nichtteiler von $xx-A$. Alle Reste von $A$ unterhalb von $A$ (ohne die Null) seien $r,r',r''$ usw., alle Nichtreste dagegen $n,n',n''$ usw. Dann ist jede Primzahl in einer der Formen $Ak+r$, $Ak+r'$, $Ak+r''$ usw. ein Teiler von $xx-A$, jede Primzahl in einer der Formen $Ak+n$, $Ak+n'$ usw. dagegen ein Nichtteiler, wobei $k$ eine unbestimmte ganze Zahl bezeichnet. Die ersten Formen nennen wir \textit{Teilerformen von $xx-A$}, die anderen \textit{Nichtteilerformen}. Von beiden Arten gibt es jeweils $\frac{1}{2}(A-1)$. Ist ferner $B$ eine zusammengesetzte ungerade Zahl und gilt $ARB$, so liegen alle Primfaktoren von $B$ in einer der \SourceLeaf{119}ersten Formen und damit auch $B$. Also ist \textit{jede} ungerade Zahl in einer Nichtteilerform ein Nichtteiler der Form $xx-A$. Dieser Satz lässt sich aber nicht umkehren: Ist $B$ ein zusammengesetzter ungerader Nichtteiler der Form $xx-A$, so sind einige Primfaktoren von $B$ Nichtteiler. Ist ihre Anzahl \textit{gerade}, so findet sich $B$ trotzdem in einer Teilerform. Siehe Artikel 99.

\textit{Beispiel.} So erhält man für $A=-11$ die folgenden Teilerformen von $xx+11$: $11k+\nobreak{}1,\allowbreak{}3,\allowbreak{}4,\allowbreak{}5,\allowbreak{}9$; die Nichtteilerformen sind dagegen $11k+\nobreak{}2,\allowbreak{}6,\allowbreak{}7,\allowbreak{}8,\allowbreak{}10$. Also ist $-11$ ein Nichtrest aller ungeraden Zahlen in einer der letzteren Formen, ein Rest dagegen aller Primzahlen, die zu einer der ersteren gehören.

Ähnliche Formen gibt es für die Teiler und Nichtteiler von $xx-A$, welche Zahl $A$ auch bezeichnet. Offensichtlich brauchen wir aber nur die Werte von $A$ zu betrachten, die durch kein Quadrat teilbar sind. Ist nämlich $A=a^2A'$, so sind alle Teiler\HistoricalNote{*)}{Gemeint sind die zu $A$ teilerfremden.} von $xx-A$ auch Teiler von $xx-A'$; dasselbe gilt für die Nichtteiler. --- Wir unterscheiden drei Fälle: 1) wenn $A$ die Form $+(4n+1)$ oder $-(4n-1)$ hat; 2) wenn $A$ die Form $-(4n+1)$ oder $+(4n-1)$ hat; 3) wenn $A$ gerade ist, also die Form $\pm(4n+2)$ hat.

\Needspace{10\baselineskip}
\Article{148}
\textit{Erster Fall:} wenn $A$ die Form $+(4n+1)$ oder $-(4n-1)$ hat. Man zerlege $A$ in seine Primfaktoren und gebe denjenigen der Form $4n+1$ ein positives, denen der Form $4n-1$ ein negatives Vorzeichen (sodass ihr Produkt $=A$ wird). Diese Faktoren seien $a,b,c,d$ usw. Man teile alle Zahlen unterhalb von $A$, die zu $A$ teilerfremd sind, in zwei Klassen ein. In die erste kommen alle Zahlen, die von keiner der Zahlen $a,b,c,d$ usw. ein Nichtrest sind oder von zweien oder vieren oder allgemein von einer geraden Anzahl; in die zweite diejenigen, die von einer der Zahlen $a,b,c$ usw. ein Nichtrest sind oder von dreien usw. oder allgemein von einer ungeraden Anzahl. Die ersten bezeichne man mit $r,r',r''$ usw., die zweiten mit $n,n',n''$ usw. Dann sind $Ak+r$, $Ak+r'$, $Ak+r''$ usw. Teilerformen von $xx-A$, dagegen $Ak+n$, $Ak+n'$ usw. Nichtteilerformen von $xx-A$ (\textit{das heißt: Jede Primzahl außer $2$ ist ein Teiler oder Nichtteiler von $xx-A$, je nachdem, ob sie in einer der ersten oder zweiten Formen liegt}). Ist nämlich $p$ eine \SourceLeaf{120}positive Primzahl und ein Rest oder Nichtrest einer der Zahlen $a,b,c$ usw., so ist diese Zahl selbst ein Rest oder Nichtrest von $p$ (nach dem Fundamentalsatz). Gibt es also unter $a,b,c$ usw. $m$ Zahlen, von denen $p$ ein Nichtrest ist, so sind ebenso viele Nichtreste von $p$. Liegt demnach $p$ in einer der ersten Formen, so ist $m$ gerade und es gilt $ARp$; liegt sie in einer der zweiten, so ist $m$ ungerade und es gilt $ANp$.

\textit{Beispiel.} Sei $A=+105=-3\times+5\times-7$. Dann sind $r,r',r''$ usw. die folgenden Zahlen: $1,\allowbreak{}4,\allowbreak{}16,\allowbreak{}46,\allowbreak{}64,\allowbreak{}79$ (die von keiner der Zahlen $3,\allowbreak{}5,\allowbreak{}7$ Nichtreste sind); $2,\allowbreak{}8,\allowbreak{}23,\allowbreak{}32,\allowbreak{}53,\allowbreak{}92$ (Nichtreste von $3,\allowbreak{}5$); $26,\allowbreak{}41,\allowbreak{}59,\allowbreak{}89,\allowbreak{}101,\allowbreak{}104$ (Nichtreste von $3,\allowbreak{}7$); $13,\allowbreak{}52,\allowbreak{}73,\allowbreak{}82,\allowbreak{}97,\allowbreak{}103$ (Nichtreste von $5,\allowbreak{}7$). --- Die Zahlen $n,n',n''$ usw. sind dagegen $11,\allowbreak{}29,\allowbreak{}44,\allowbreak{}71,\allowbreak{}74,\allowbreak{}86$; $22,\allowbreak{}37,\allowbreak{}43,\allowbreak{}58,\allowbreak{}67,\allowbreak{}88$; $19,\allowbreak{}31,\allowbreak{}34,\allowbreak{}61,\allowbreak{}76,\allowbreak{}94$; $17,\allowbreak{}38.47,\allowbreak{}62,\allowbreak{}68,\allowbreak{}83$. Die ersten sechs sind Nichtreste von $3$, die nächsten sechs Nichtreste von $5$, dann folgen die Nichtreste von $7$ und schließlich die Nichtreste aller drei zugleich.

Aus der Kombinationslehre und den Artikeln 32, 96 folgt leicht, dass die Anzahl der Zahlen $r,r',r''$ usw.
\[
=t\left(1+\frac{l\, .\, l-1}{1\, .\, 2}+\frac{l\, .\, l-1\, .\, l-2\, .\, l-3}{1\, .\, 2\, .\, 3\, .\, 4}+ .\, .\, .\, .\, .\right)
\]
und die Anzahl der Zahlen $n,n',n''$ usw.
\[
=t\left(l+\frac{l\, .\, l-1\, .\, l-2}{1\, .\, 2\, .\, 3}+\frac{l\, .\, l-1\, .\, .\, .\, .\, l-4}{1\, .\, 2\, .\, .\, .\, .\, 5}+ .\, .\, .\right)
\]
ist, wobei $l$ die Anzahl der Zahlen $a,b,c$ usw. bezeichnet;
\[
t=2^{-l}(a-1)(b-1)(c-1)\text{ usw.}
\]
und jede der beiden Reihen ist fortzusetzen, bis sie abbricht. (Es gibt nämlich $t$ Zahlen, die Reste aller $a,b,c$ usw. sind, $\frac{t\, .\, l\, .\, l-1}{1\, .\, 2}$, die Nichtreste von zweien sind, usw.; der Kürze wegen können wir diesen Beweis aber nicht ausführlicher entwickeln.) Die Summe jeder Reihe\HistoricalNote{*)}{Ohne den Faktor $t$.} ist $=2^{l-1}$. Die erste entsteht aus
\[
1+(l-1)+\frac{l-1\, .\, l-2}{1\, .\, 2}+ .\, .\, .
\]
durch Zusammennehmen des zweiten und dritten, des vierten und fünften Gliedes usw.; die zweite aus derselben Reihe durch Zusammennehmen des ersten und zweiten, des dritten und vierten Gliedes usw. Also gibt es ebenso viele Teilerformen von $xx-A$ wie Nichtteilerformen, nämlich $\frac{1}{2}(a-1)(b-1)(c-1)$ usw.

\SourceLeaf{121}\Needspace{10\baselineskip}
\Article{149}
Den \textit{zweiten und dritten Fall} können wir hier gemeinsam betrachten. Man kann hier stets $A$ $=(-1)Q$ oder $=(+2)Q$ oder $=(-2)Q$ setzen, wobei $Q$ eine Zahl der Form $+(4n+1)$ oder $-(4n-1)$ bezeichnet, wie wir sie im vorigen Artikel betrachtet haben. Allgemein sei $A=\alpha Q$, wobei $\alpha$ entweder $=-1$ oder $=\pm2$ ist. Dann ist $A$ ein Rest aller Zahlen, von denen entweder beide $\alpha$ und $Q$ Reste sind oder keine von beiden, dagegen ein Nichtrest aller Zahlen, von denen genau eine der Zahlen $\alpha,Q$ ein Nichtrest ist. Hieraus ergeben sich die Teiler- und Nichtteilerformen von $xx-A$ leicht. Ist $\alpha=-1$, so teile man alle Zahlen unterhalb von $4A$, die dazu teilerfremd sind, in zwei Klassen. In die erste kommen diejenigen, die in einer Teilerform von $xx-Q$ und zugleich in der Form $4n+1$ liegen, und diejenigen, die in einer Nichtteilerform von $xx-Q$ und zugleich in der Form $4n+3$ liegen; in die zweite die übrigen. Die ersten seien $r,r',r''$ usw., die zweiten $n,n',n''$ usw. Dann ist $A$ ein Rest aller Primzahlen in einer der Formen $4Ak+r$, $4Ak+r'$, $4Ak+r''$ usw., ein Nichtrest dagegen aller Primzahlen in einer der Formen $4Ak+n$, $4Ak+n'$ usw. --- Ist $\alpha=\pm2$, so teile man alle Zahlen unterhalb von $8Q$, die dazu teilerfremd sind, in zwei Klassen. In die erste kommen diejenigen, die in einer Teilerform von $xx-Q$ und zugleich in einer der Formen $8n+1$, $8n+7$ beim oberen Zeichen oder $8n+1$, $8n+3$ beim unteren liegen; außerdem diejenigen, die in einer Nichtteilerform von $xx-Q$ und zugleich in einer der Formen $8n+3$, $8n+5$ beim oberen Zeichen oder $8n+5$, $8n+7$ beim unteren liegen. In die zweite kommen die übrigen. Bezeichnet man die Zahlen der ersten Klasse mit $r,r',r''$ usw. und die der zweiten mit $n,n',n''$ usw., so ist $\pm2Q$ ein Rest aller Primzahlen in einer der Formen $8Qk+r$, $8Qk+r'$, $8Qk+r''$ usw., ein Nichtrest dagegen aller Primzahlen in einer der Formen $8Qk+n$, $8Qk+n'$, $8Qk+n''$ usw. Auch hier lässt sich leicht beweisen, dass es ebenso viele Teilerformen von $xx-A$ wie Nichtteilerformen gibt.

\textit{Beispiel.} So ergibt sich $+10$ als Rest aller Primzahlen in einer der Formen $40k+\nobreak{}1,\allowbreak{}3,\allowbreak{}9,\allowbreak{}13,\allowbreak{}27,\allowbreak{}31,\allowbreak{}37,\allowbreak{}39$, als Nichtrest dagegen aller Primzahlen in einer der Formen $40k+\nobreak{}7,\allowbreak{}11,\allowbreak{}17,\allowbreak{}19,\allowbreak{}21,\allowbreak{}23,\allowbreak{}29,\allowbreak{}33$.

\SourceLeaf{122}\Needspace{8\baselineskip}
\Article{150}
Diese Formen haben mehrere recht bemerkenswerte Eigenschaften; wir führen aber nur eine an. Ist $B$ eine zusammengesetzte, zu $A$ teilerfremde Zahl, unter deren Primfaktoren $2m$ in einer Nichtteilerform von $xx-A$ liegen, so liegt $B$ in einer Teilerform von $xx-A$. Ist dagegen die Anzahl der Primfaktoren von $B$, die in einer Nichtteilerform von $xx-A$ liegen, ungerade, so liegt auch $B$ in einer Nichtteilerform. Den Beweis, der nicht schwierig ist, lassen wir weg. Daraus folgt, dass nicht nur jede Primzahl, sondern auch jede zusammengesetzte ungerade, zu $A$ teilerfremde Zahl in einer Nichtteilerform ein Nichtteiler ist: Ein Primfaktor einer solchen Zahl muss notwendig ein Nichtteiler sein.

\Needspace{12\baselineskip}
\begin{center}\textit{Über die Arbeiten anderer zu diesen Untersuchungen.}\end{center}

\Article{151}
Der Fundamentalsatz, der gewiss zu den elegantesten Ergebnissen auf diesem Gebiet gehört, wurde bisher von niemandem in der einfachen Form ausgesprochen, in der wir ihn oben angegeben haben. Das ist umso erstaunlicher, als dem berühmten Euler bereits gewisse andere darauf beruhende Aussagen bekannt waren, von denen man leicht zu ihm hätte zurückkommen können. Er wusste, dass es bestimmte Formen gibt, die alle Primteiler der Zahlen der Form $xx-A$ enthalten, und andere, die alle primen Nichtteiler von Zahlen derselben Form umfassen und die ersten Formen ausschließen; auch hatte er ein Verfahren gefunden, diese Formen zu bestimmen. Alle seine Versuche, zu einem Beweis zu gelangen, blieben aber erfolglos und machten diese durch Beispiele gefundene Wahrheit nur wahrscheinlicher. In einer Abhandlung, \textit{Novae demonstrationes circa divisores numerorum formae $xx+nyy$}, die am 20. November 1775 in der Petersburger Akademie vorgetragen wurde und nach dem Tod dieses großen Mannes in \textit{Bd.} I der \textit{Nov. Act.} dieser Akademie, S. 47 ff., erhalten ist, scheint er allerdings geglaubt zu haben, sein Ziel erreicht zu haben. Doch hier schlich sich ein Fehler ein: Auf S. 65 nahm er \textit{stillschweigend} an, dass solche Teiler- und Nichtteilerformen existieren\HistoricalNote{*}{Dass es nämlich Zahlen $r,r',r''$ usw.; $n,n',n''$ usw. gibt, alle verschieden und $<4A$, sodass alle Primteiler von $xx-A$ in einer der Formen $4Ak+r$, $4Ak+r'$ usw. und alle primen Nichtteiler in einer der Formen $4Ak+n$, $4Ak+n'$ usw. liegen (wobei $k$ eine unbestimmte Zahl bezeichnet).}. Daraus war leicht herzuleiten, \textit{welche} Formen es sein müssen; das Verfahren, mit dem er diese Annahme begründete, scheint jedoch nicht \SourceLeaf{123}geeignet. In einer anderen Schrift, \textit{De criteriis aequationis $fxx+gyy=hzz$ utrumque resolutionem admittat necne}, Opusc. Anal., Bd. I (wobei $f,g,h$ gegeben und $x,y,z$ unbestimmt sind), fand er anhand von Beispielen: Ist die Gleichung für einen Wert $h=s$ lösbar, so ist sie auch für jeden anderen Wert lösbar, der $s$ nach dem Modul $4fg$ kongruent ist, sofern er eine Primzahl ist. Aus diesem Satz lässt sich die erwähnte Annahme ohne Schwierigkeit beweisen. Aber auch der Beweis dieses Satzes entzog sich allen seinen Bemühungen\HistoricalNote{*)}{Wie er selbst an der angegebenen Stelle, S. 216, zugibt: „Ein Beweis dieses höchst eleganten Satzes fehlt noch, nachdem viele ihn schon lange vergeblich gesucht haben.\,\ldots\ Wer einen Beweis dieses Satzes findet, wird also sehr viel geleistet haben.“ --- Wie sehr dieser unsterbliche Mann einen Beweis dieses Satzes und anderer wünschte, die nur besondere Fälle des Fundamentalsatzes sind, zeigen viele andere Stellen der Opusc. Anall. Vergleiche \textit{Additamentum ad diss.} VIII, Bd. I, und \textit{diss.} XIII, Bd. II, sowie mehrere bereits an verschiedenen Stellen angeführte Abhandlungen in den Comment. Petrop.}. Das ist nicht erstaunlich, da man unserer Ansicht nach vom Fundamentalsatz ausgehen musste. Im Übrigen wird sich die Richtigkeit dieses Satzes unmittelbar aus dem ergeben, was wir im nächsten Abschnitt behandeln.

Nach Euler beschäftigte sich der bekannte Le Gendre mit demselben Gegenstand in seiner ausgezeichneten Abhandlung \textit{Recherches d'analyse indéterminée}, Hist. de l'Ac. des Sc., 1785, S. 465 ff. Dort gelangte er zu einem Satz, der inhaltlich mit dem Fundamentalsatz übereinstimmt: Bezeichnen $p,q$ zwei positive Primzahlen, so sind die im Betrag kleinsten Reste der Potenzen $p^{\frac{q-1}{2}}$, $q^{\frac{p-1}{2}}$ nach den Moduln $q,p$ entweder beide $+1$ oder beide $-1$, wenn $p$ oder $q$ die Form $4n+1$ hat. Haben dagegen sowohl $p$ als auch $q$ die Form $4n+3$, so ist der eine kleinste Rest $+1$, der andere $-1$ (S. 516). Daraus folgt nach Artikel 106, dass die \textit{Beziehung} (im Sinne von Artikel 146) von $p$ zu $q$ und von $q$ zu $p$ \textit{dieselbe} ist, wenn $p$ oder $q$ die Form $4n+1$ hat, dagegen \textit{entgegengesetzt}, wenn sowohl $p$ als auch $q$ die Form $4n+3$ haben. Dieser Satz ist unter den Aussagen aus Artikel 131 enthalten; er folgt auch aus 1, 3, 9 des Artikels 133. Umgekehrt lässt sich der Fundamentalsatz aus ihm herleiten. Le Gendre versuchte auch einen Beweis, über den wir, da er äußerst einfallsreich ist, im nächsten Abschnitt ausführlicher sprechen werden. Da er aber darin einiges ohne Beweis voraussetzte (wie er selbst auf S. 520 zugibt: \textit{Wir haben nur angenommen} usw.), was teils bisher von niemandem bewiesen wurde, teils jedenfalls unserer Ansicht nach ohne den Fundamentalsatz selbst nicht bewiesen werden kann, scheint sein eingeschlagener Weg nicht zum Ziel führen zu können; unser Beweis muss daher als der erste gelten. --- Weiter unten werden wir noch \textit{zwei andere Beweise} dieses äußerst wichtigen Satzes geben, die vom vorigen und voneinander völlig verschieden sind.

\SourceLeaf{124}\Needspace{38\baselineskip}
\begin{center}\textit{Quadratische Kongruenzen mit linearem Glied.}\end{center}

\Article{152}
Bisher haben wir die reine Kongruenz $xx\equiv A\ (\mathrm{mod.}\,m)$ behandelt und gezeigt, wie sich ihre Lösbarkeit entscheiden lässt. Die Ermittlung \textit{der Wurzeln selbst} wurde in Artikel 105 auf den Fall zurückgeführt, dass $m$ eine Primzahl oder Primzahlpotenz ist; der zweite Fall in Artikel 101 auf den, dass $m$ eine Primzahl ist. Für diesen Fall umfasst das in Artikel 61 ff. Angegebene zusammen mit dem, was wir in den Abschnitten V und VIII behandeln werden, fast alles, was sich durch direkte Methoden gewinnen lässt. Diese sind aber dort, wo sie anwendbar sind, meist unendlich viel umständlicher als die indirekten Methoden, die wir in Abschnitt VI behandeln werden, und daher eher wegen ihrer Schönheit als wegen ihres praktischen Nutzens bemerkenswert. --- \textit{Quadratische Kongruenzen mit linearem Glied} lassen sich leicht auf reine zurückführen. Soll die gegebene Kongruenz
\[
axx+bx+c\equiv 0
\]
nach dem Modul $m$ gelöst werden, so ist dazu die folgende Kongruenz äquivalent:
\[
4aaxx+4abx+4ac\equiv 0\ (\mathrm{mod.}\,4am)
\]
\textit{Das heißt,} jede Zahl, die einer von beiden genügt, genügt auch der anderen. Diese lässt sich so schreiben:
\[
(2ax+b)^2\equiv bb-4ac\ (\mathrm{mod.}\,4am)
\]
Daraus lassen sich alle Werte von $2ax+b$ unterhalb von $4am$ finden, sofern es welche gibt. Bezeichnet man sie mit $r,r',r''$ usw., so erhält man alle Lösungen der gegebenen Kongruenz aus den Lösungen der Kongruenzen
\[
2ax\equiv r-b,\quad 2ax\equiv r'-b\text{ usw.}\ (\mathrm{mod.}\,4am)
\]
die wir in Abschnitt II finden lehrten. Außerdem lässt sich die Lösung meist durch verschiedene Kunstgriffe verkürzen. \textit{Zum Beispiel} kann man anstelle der gegebenen eine äquivalente Kongruenz
\[
a'xx+2b'x+c'\equiv 0
\]
finden, in der $a'$ die Zahl $m$ teilt. Der Kürze wegen können wir dies hier nicht erklären; dazu kann man den letzten Abschnitt heranziehen.

\par\endgroup
